\chapter{각운동량, 스핀, 통계}

\section{고전적 회전에서 고유 스핀까지}

형식적으로 다루는 주제는 각운동량이지만, 실제로 살펴볼 대상은 스핀이다. 스핀은 기본 입자의 가장 단순하면서도 직관에 어긋나는 성질 가운데 하나이며, 그 구조는 각운동량의 수학을 따른다.

각운동량은 벡터다. 방향은 회전축을 따라 정해진다. 그 축을 따라 가능한 방향은 두 가지이므로, 어느 쪽을 택할지 정하는 약속이 필요하다. 회전 방향으로 오른손 손가락을 감아 쥐면 엄지가 각운동량 벡터의 방향을 가리킨다. 이는 정리가 아니라 정의다.

각운동량의 크기는 질량이 얼마나 있는지, 질량이 어떻게 분포하는지, 계의 크기가 얼마인지, 각속도가 얼마인지에 따라 달라진다. 질량 분포와 크기는 관성 모멘트에 모아 표현한다. 따라서 큰 물체와 작은 물체는 회전 속도와 에너지가 크게 달라도 같은 각운동량을 가질 수 있다.

보통의 복합 물체는 개념적으로 서로 다른 두 종류의 각운동량을 가질 수 있다. 컵이 외부 축을 중심으로 원운동하면, 컵 자체를 관통하는 어떤 축을 중심으로 컵이 돌지 않더라도 질량 중심이 그 축 주위를 움직이므로 각운동량을 갖는다. 이것이 궤도 각운동량이다. 스핀은 내부 축을 중심으로 한 회전을 뜻한다.

깔끔한 정의는 질량중심계에서 내릴 수 있다. 질량중심이 정지해 있어 전체 보통 운동량이 사라지는 기준계로 가자. 그 기준계에서 남는 각운동량은 모두 스핀 각운동량이다. 물체 내부까지 정지해 있을 필요는 없다. 질량중심만 정지해 있으면 된다.

농구공에서는 이 구분이 어렵지 않다. 정지한 농구공과 회전시킨 같은 농구공은 서로 다른 두 상태에 있는 동일한 물체임이 분명하다. 양자역학에서는 각운동량이 불연속이기 때문에 이 그림이 복잡해진다. 각운동량이 0인 상태와 처음으로 회전하는 상태 사이에는 허용된 각운동량이 연속적으로 이어져 있지 않다. 그러므로 처음 회전하는 상태가 움직이는 동일한 물체인지, 아니면 완전히 다른 물체인지 물을 수 있다.

표현의 일부는 관례에 따른 것이지만, 물리적 질문은 명확하다. 처음으로 0이 아닌 회전 상태에 도달하려면 에너지가 얼마나 필요한가? 에너지 간격이 작으면 들뜬 물체는 여전히 쉽게 알아볼 수 있다. 간격이 매우 크면 들뜸으로 물체가 크게 달라질 수 있다. 원자에 회전 들뜸 에너지를 지나치게 많이 주면, 원자가 그대로 조용히 회전하는 대신 전자가 떨어져 나가 원자 자체가 파괴될 수 있다.

이제 알려진 고유 스핀은 잠시 제쳐 두고 전자를 생각해 보자. 전자를 내부 축을 중심으로 회전시킬 수 있을까? 수학적 점에는 회전할 가시적인 내부 구조가 없지만, 실험에서 전자의 크기를 확인할 때 관찰 가능한 작은 단단한 표면을 제공해 그것이 도는 모습을 보는 것은 아니다. 전자에 극도로 작은 유한한 크기가 있다면, 처음으로 가능한 추가 회전 들뜸에 도달하는 데 너무 많은 에너지가 필요할 수도 있다. 크기가 작아질수록 에너지가 커지므로, 필요한 에너지는 천문학적으로 클 수 있고 플랑크 에너지 정도일 수도 있다. 이 추정은 의도적으로 잠정적이다.\footnote{이 추정 바로 뒤에 나오는 짧은 대목은 추가적인 실질적 재구성을 뒷받침할 만큼 충분히 명확하지 않다.}

따라서 기본 입자의 스핀 크기는 고정된 특성으로 취급한다. 전자는 점점 더 큰 고유 스핀 크기의 범위를 갖지 않는다. 다른 점에서는 전자와 같은 들뜸이라도 전체 스핀이 다르면 다른 입자로 분류된다.

고정된 크기가 방향까지 고정하는 것은 아니다. 각운동량의 크기로 부분적으로 식별되는 \emph{스핀트론}이라는 입자를 생각해 보자. 스핀이 한 축을 따라 향한다면, 회전 대칭성에 따라 같은 종류의 입자에서 스핀이 다른 축을 따라 향하는 경우도 가능해야 한다. 어느 공간 방향도 본질적으로 특별하지 않다. 스핀의 양은 고정된 채 방향은 바뀔 수 있다. 따라서 이 문제에는 서로 모순되어 보이는 두 요소, 즉 양자화된 크기와 자유롭게 회전할 수 있는 방향이 함께 있다. 각운동량 대수는 이 둘이 어떻게 양립하는지 보여 줄 것이다.

\begin{classroomqa}\lecturetimestamp{00:12:25}\audiencequestion{작은 입자를 회전시키는 데 왜 에너지가 필요한가?}\lecturerresponse{질량이 \(M\)이고 특성 반지름이 \(R\)인 물체의 관성 모멘트는 대략 다음과 같다.\begin{equation}
I\sim cMR^2,
\end{equation}여기서 수치 인자 \(c\)는 물체의 모양과 질량 분포에 따라 달라진다. 바깥쪽에 있는 물질이 속력 \(v\)로 움직인다면, 그 보통 운동량은 대략 \(Mv\)이고 각운동량은 운동량에 거리를 곱한 값이다.\begin{equation}
L\sim MvR.
\end{equation}같은 내용을 \(L=I\omega\), \(v\sim\omega R\)를 사용해 쓸 수도 있다. 회전 운동 에너지는 다음과 같다.\begin{equation}
E_{\mathrm{rot}}=\frac{L^2}{2I}.
\end{equation}\(M\)와 \(L\)가 고정되어 있으면,\begin{equation}
E_{\mathrm{rot}}\sim\frac{L^2}{2cMR^2}\propto\frac{1}{R^2}.
\end{equation}물체가 작을수록 주어진 각운동량에 더 많은 에너지가 든다. 이 역제곱 의존성은 고전적인 결과다. 여기에 더해지는 양자역학적 사실은 각운동량이 불연속적인 단계로 변하므로, 처음 허용되는 변화량을 임의로 작게 할 수 없다는 것이다.}\end{classroomqa}

이제 사용할 수학은 추상적이고 처음에는 다소 직관에 어긋난다. 하지만 소수의 대수적 규칙만으로 실험으로 검증할 수 있는 사실을 이끌어 낼 수 있다는 점에서 흥미롭다.

\section{궤도 각운동량}

원점을 중심으로 원 궤도를 도는 입자 하나부터 생각해 보자. 입자의 각운동량 크기는 운동량에 원점에서 수직으로 잰 거리를 곱한 값이다. 평면에서는 한 방향의 회전을 양으로, 반대 방향의 회전을 음으로 정할 수 있다. 이 부호는 선택한 축을 기준으로 한 각운동량 벡터의 방향을 나타낸다.

이것은 아직 스핀이 아니다. 입자가 보통의 운동량을 가지기 때문에 각운동량도 갖는다. 이제 지름의 양 끝에 입자 두 개를 놓고 같은 방향으로 회전하게 하자. 보통의 운동량은 서로 상쇄되어 질량중심은 정지하지만, 각운동량은 같은 방향을 향하므로 서로 더해진다. 이 계의 전체 운동량은 0이지만 전체 각운동량은 0이 아니다. 스핀의 고전적 모형이다.

원점에 대한 입자의 위치를 다음과 같이 두고\begin{equation}
\vec r=(x,y,z)=(x_1,x_2,x_3),
\end{equation}운동량을 다음과 같이 두자.\begin{equation}
\vec p=(p_x,p_y,p_z)=(p_1,p_2,p_3).
\end{equation}각운동량은 벡터곱으로 주어진다.\begin{equation}
\vec L=\vec r\times\vec p.
\end{equation}인자의 순서를 바꾸면 벡터의 방향도 뒤집힌다. \(\vec p\times\vec r=-\vec r\times\vec p\). 표준 약속은 \(\vec r\times\vec p\)이며, 그 방향은 오른손 법칙으로 정한다.\footnote{\raggedright 강의 영상 프레임 \nolinkurl{lecture_07_figure_02.png}, 00:18:49.\par}

\begin{figure}[h]\centering
\includegraphics[width=0.72\textwidth]{lecture_07_figure_02.png}\par
\begin{tikzpicture}[scale=1.0]
\draw[->] (-0.6,0) -- (4.8,0) node[right] {$x$};
\draw[->] (0,-0.6) -- (0,3.8) node[above] {$y$};
\coordinate (P) at (3.2,2.3);
\draw[->,thick] (0,0) -- (P) node[midway,above left] {$\vec r$};
\draw[dashed] (P) -- (3.2,0);
\draw[dashed] (P) -- (0,2.3);
\draw[->,thick] (P) -- ++(-1.0,1.25) node[above] {$\vec p$};
\node at (4.15,2.75) {$\vec L=\vec r\times\vec p$};
\end{tikzpicture}\caption{위치가 \(\vec r\)이고 운동량이 \(\vec p\)인 입자에서 각운동량 \(\vec L=\vec r\times\vec p\)은 두 벡터가 놓인 평면에 수직이다.}\lectureframe{00:18:49}\end{figure}

데카르트 성분으로 쓰면 다음과 같다.\begin{equation}
\begin{aligned}
L_x&=yp_z-zp_y,\qquad
L_y=zp_x-xp_z,\\
L_z&=xp_y-yp_x.
\end{aligned}
\end{equation}첫 번째 식만 기억하고 첨자 \[
x\longrightarrow y\longrightarrow z\longrightarrow x.
\]를 순환시키면 충분하다. 성분의 첨자와 우변의 첨자를 함께 순환하면 상대적인 부호를 바꾸지 않고 다음 식을 얻는다. 이 기억법이 유용한 이유는 부호를 실수로 뒤집기 쉽기 때문이다.

점입자의 경우 \(\vec r\times\vec p\)는 선택한 원점에 대한 궤도 각운동량이다. 복합계에서는 모든 구성 입자의 각운동량을 벡터로 더한다.
\begin{equation}
\vec L_{\mathrm{tot}}=\sum_a \vec r_a\times\vec p_a.
\end{equation}
이 합을 질량중심계에서 계산하면 다음과 같이 전체 운동량이 상쇄되어도 각운동량은 0이 아닐 수 있다.
\[
\sum_a\vec p_a=0.
\]
이렇게 남는 각운동량이 복합계의 스핀이다.

\section{양자 연산자와 정준 교환자}

양자역학에서는 관측 가능한 물리량을 연산자로 나타낸다. 따라서 위치, 운동량, 각운동량의 각 성분은 연산자다. 여기서는 이들을 미분 연산자로 구체적으로 나타내는 방법까지 알 필요는 없다. 교환 관계만으로도 각운동량 스펙트럼을 유도하기에 충분하다.

두 연산자 \(A\)와 \(B\)에 대해 다음과 같이 정의한다.\begin{equation}
[A,B]=AB-BA.
\end{equation}두 관측량이 교환하면 두 값을 동시에 지정할 수 있다. 교환하지 않으면 양자역학에 따라 두 값을 동시에 정확히 지정할 수 없다.

위치의 세 좌표 성분은 함께 측정할 수 있으므로 서로 교환한다.\begin{equation}
[x_i,x_j]=0.
\end{equation}특히 모든 연산자는 자기 자신과 교환한다. 어떤 관측량이 그 관측량 자체를 다시 측정하는 것과 양립하지 않는다고 보는 것은 말이 되지 않는다.

운동량 성분들도 마찬가지로 서로 교환한다.\begin{equation}
[p_i,p_j]=0.
\end{equation}위치 성분은 다른 방향의 운동량 성분과 교환하지만, 같은 방향의 운동량 성분과는 교환하지 않는다. 세 가지 정준 관계를 함께 쓰면 다음과 같다.\begin{equation}
[x_i,x_j]=0,\qquad
[p_i,p_j]=0,\qquad
[x_i,p_j]=i\hbar\,\delta_{ij}.
\end{equation}따라서 예를 들면 다음과 같다.\[
[x,p_y]=0,\qquad [x,p_x]=i\hbar,
\]연산자 순서를 바꾸면 부호가 바뀐다.\begin{equation}
[p_j,x_i]=-i\hbar\,\delta_{ij}.
\end{equation}

각운동량과 플랑크 상수는 같은 차원을 갖는다. 각운동량은 거리와 운동량의 곱이고, 불확정성 관계의 차원은 다음과 같다.\begin{equation}
\Delta x\,\Delta p\sim\hbar
\end{equation}정확한 부등식에 익숙한 수치 인자를 제외한 형태다. 그러므로 각운동량이 \(\hbar\) 단위로 양자화되리라고 보는 것은 자연스럽지만, 아직 양자화를 유도한 것은 아니다.

\section{각운동량 대수의 유도}

먼저 각운동량의 두 성분 사이 교환자를 계산하자.\[
L_x=yp_z-zp_y,\qquad
L_y=zp_x-xp_z,
\]다음과 같이 된다.\begin{align}
[L_x,L_y]
&=[yp_z-zp_y,\;zp_x-xp_z] \\
&=[yp_z,zp_x]-[yp_z,xp_z] \notag\\
&\quad{}-[zp_y,zp_x]+[zp_y,xp_z].
\end{align}

\begin{classroomqa}\lecturetimestamp{00:31:41}\audiencequestion{위에서 아래로 세 식을 순환시키는 대신, 기호를 왼쪽에서 오른쪽으로 순환시켜 성분 식을 기억할 수 있나요?}\lecturerresponse{네. 어느 쪽 순환 기억법을 써도 된다. 기억법의 목적은 부호를 추적하면서 성분을 올바른 순환 순서로 만드는 것뿐이다.}\end{classroomqa}

교환자는 한 곱의 인자들을 다른 곱의 인자들 사이로 옮길 수 있는지 살펴보면 이해할 수 있다. 다음 식에서
\[
[yp_z,xp_z],
\]
연산자 \(y\)는 \(x\) 및 \(p_z\)와 각각 가환하고, \(p_z\)는 \(x\) 및 자기 자신과 가환한다. 따라서 서로 방해하는 항이 없어 이 항은 0이 된다. 같은 논리로 다음을 얻는다.
\[
[zp_y,zp_x]=0.
\]
곱의 쌍 가운데 두 쌍만 남는다.

첫 번째 쌍에서는 \(z\)을 통과하는 \(p_z\)에서만 걸린다.\begin{align}
[yp_z,zp_x]
&=yp_zzp_x-zp_xyp_z \\
&=y(p_zz-zp_z)p_x \\
&=y[p_z,z]p_x \\
&=-i\hbar\,yp_x.
\end{align}정준 교환자의 순서는 중요하다.\[
[p_z,z]=-[z,p_z]=-i\hbar.
\]

남은 다른 쌍에서는 \(L_x\)와 \(L_y\)의 정의에 들어 있는 음수 부호 두 개가 곱해져 양수가 된다. \(p_z\)을 통과하는 \(z\)에서만 걸린다.\begin{align}
[zp_y,xp_z]
&=zp_yxp_z-xp_zzp_y \\
&=x(zp_z-p_zz)p_y \\
&=x[z,p_z]p_y \\
&=i\hbar\,xp_y.
\end{align}두 기여를 합하면 다음과 같다.\begin{align}
[L_x,L_y]
&=-i\hbar\,yp_x+i\hbar\,xp_y \\
&=i\hbar(xp_y-yp_x) \\
&=i\hbar L_z.
\end{align}

이것이 중요한 단순화다. 각운동량의 두 성분을 교환하면 세 번째 성분이 나오고, 위치나 운동량 연산자는 남지 않는다. 순환 순열을 적용하면 다음을 얻는다.\begin{equation}
\begin{aligned}
[L_x,L_y]&=i\hbar L_z,\qquad
[L_y,L_z]=i\hbar L_x,\\
[L_z,L_x]&=i\hbar L_y.
\end{aligned}
\end{equation}같은 내용을 다음과 같이 쓸 수도 있다.\begin{equation}
[L_i,L_j]=i\hbar\,\epsilon_{ijk}L_k.
\end{equation}

대수는 각운동량 연산자들만으로 닫혀 있다. 이 관계들을 얻고 나면, 유도에 사용한 \(x_i\)와 \(p_i\)는 더 이상 고려하지 않아도 된다. 같은 대수는 많은 구성 요소에 대한 합인 전체 각운동량에도 적용된다. 특히 남은 각운동량이 스핀인 질량 중심 좌표계에도 적용된다. 따라서 스핀을 작은 구성 요소들의 운동으로 먼저 구성하지 않아도 각운동량 교환 관계를 추상적으로 사용할 수 있다.

\(x\), \(y\), \(z\)축을 어떻게 택하든 상관없다. 서로 수직인 다른 관성 데카르트 좌표축으로 회전해도 프라임이 붙은 성분들은 같은 관계를 만족한다.\begin{equation}
[L'_x,L'_y]=i\hbar L'_z,
\end{equation}나머지 관계도 순환 순열로 얻는다. 대수는 회전 대칭이므로, 관성 좌표축을 어떤 것으로 골라 계산했는지는 물리적 예측에 영향을 줄 수 없다.

\begin{classroomqa}\lecturetimestamp{00:39:09}\audiencequestion{스핀 각운동량은 좌표축의 회전뿐 아니라 평행이동에도 무관한가요?}\lecturerresponse{네. 고유 스핀은 평행이동에 무관하다. 원점에 대한 궤도 각운동량은 원점을 평행이동하면 달라질 수 있지만, 질량 중심 좌표계에서 남는 스핀은 그런 평행이동에 의존하지 않는다.}\end{classroomqa}

추상 대수로는 끝없이 많은 항등식을 만들 수 있지만, 대부분은 물리적으로 흥미롭지 않다. 항등식이 측정 가능한 사실이 되거나 측정의 조작적 의미를 드러낼 때 비로소 가치가 생긴다. 지금까지의 전개는 여전히 주로 대수적이다. 다음 단계에서는 스핀의 불연속 구조를 이끌어 낸다.

\begin{classroomqa}\lecturetimestamp{00:40:42}\audiencequestion{선택한 \(x\), \(y\), \(z\)축이 입자들과 함께 회전한다면 어떻게 되나요?}\lecturerresponse{여기서는 관성계만 다룬다. 회전 좌표계는 비관성계이며 원심력과 코리올리 힘 같은 효과를 도입한다. 각운동량은 공동 회전 좌표계로 바꾸어도 불변이 아니다. 관성계에서 명백히 회전하는 계가 그 계와 함께 회전하는 좌표계에서는 정지해 보일 수 있다. 여기서 사용하는 회전 불변성은 서로 다른 방향으로 놓인 관성 데카르트 좌표축에 관한 것이지, 시간에 따라 회전하는 축에 관한 것이 아니다.}\end{classroomqa}

\section{사다리 연산자와 양자화 축}

다음 두 조합을 도입하자.\begin{equation}
L_+=L_x+iL_y,
\qquad
L_-=L_x-iL_y.
\end{equation}이들은 조화 진동자의 생성 연산자와 소멸 연산자에 대응한다. 여기서는 선택한 각운동량 성분을 올리거나 내리는 역할을 한다.

\(z\)축이 특별한 것은 아니지만, 어떤 성분을 지정하려면 축 하나를 택해야 한다. \(z\)를 양자화 축으로 정하고 다음과 같이 쓰자.\begin{equation}
L_z\lvert m\rangle=m\hbar\,\lvert m\rangle.
\end{equation}그러므로 \(m\)는 \(z\) 성분을 \(\hbar\) 단위로 나타낸 측정값이다. 대수 계산 중에는 \(\hbar=1\)로 두는 것이 편리한 경우가 많지만, 여기서는 그 값을 명시적으로 유지하겠다.

교환자 \([L_+,L_-]\)는 \(L_z\)에 비례하지만, 지금 바로 필요한 것은 아니다. 먼저 중요한 것은 각 사다리 연산자와 \(L_z\) 사이의 교환자다.\[
[L_x,L_z]=-i\hbar L_y,
\qquad
[L_y,L_z]=i\hbar L_x,
\]에서 다음을 얻는다.\begin{align}
[L_+,L_z]
&=[L_x+iL_y,L_z] \\
&=[L_x,L_z]+i[L_y,L_z] \\
&=-i\hbar L_y+i(i\hbar L_x) \\
&=-\hbar(L_x+iL_y) \\
&=-\hbar L_+.
\end{align}\(iL_y\) 항의 부호를 바꾸면 다음과 같다.\begin{equation}
[L_-,L_z]=+\hbar L_-.
\end{equation}각 교환자의 순서를 뒤집어 쓰면 다음과 같은 등가 표현을 얻는다.\begin{equation}
[L_z,L_+]=\hbar L_+,
\qquad
[L_z,L_-]=-\hbar L_-.
\end{equation}

각운동량 성분 중 한 번에 하나만 정확히 정할 수 있다. 서로 다른 두 성분은 교환하지 않으므로 동시에 지정할 수 없다. 정확히 정한 성분이 양자화 축을 정한다. \(z\)를 택하는 것은 편의일 뿐, 물리적으로 특별한 선택은 아니다.

이제 \(L_+\)를 \(L_z\)의 고유상태에 작용시켜 보자. 교환자를 다음과 같이 다시 쓸 수 있다.
\begin{equation}
L_zL_+=L_+L_z+\hbar L_+.
\end{equation}
따라서
\begin{align}
L_z\bigl(L_+\lvert m\rangle\bigr)
&=\bigl(L_+L_z+\hbar L_+\bigr)\lvert m\rangle \\
&=m\hbar L_+\lvert m\rangle+\hbar L_+\lvert m\rangle \\
&=(m+1)\hbar\,L_+\lvert m\rangle.
\end{align}
\(L_+\lvert m\rangle\)가 0이 아니라면, 이는 \(L_z\)의 고유상태이며 고윳값은 \(\hbar\)만큼 증가한다. 마찬가지로
\begin{equation}
L_z\bigl(L_-\lvert m\rangle\bigr)
=(m-1)\hbar\,L_-\lvert m\rangle.
\end{equation}
따라서 \(L_-\)는 \(m\)를 한 단위 낮춘다.

여기서 이미 자명하지 않은 결론을 얻었다. \(L_z\)의 스펙트럼이 무엇으로 밝혀지든, 이웃한 값들은 정수만큼 차이 난다.\begin{equation}
\ldots,\ m-2,\ m-1,\ m,\ m+1,\ m+2,\ \ldots
\end{equation}양 끝값은 아직 정해지지 않았다.

\section{유한한 다중항과 스핀 표지}

정해진 종류의 입자는 각운동량 크기가 고정되어 있다. 각운동량의 방향은 바꿀 수 있지만 크기를 한없이 키우면 결국 다른 물체가 된다. 따라서 그 입자가 가질 수 있는 \(z\) 성분에는 최댓값이 있어야 한다.

올림 연산자가 있는데도 사다리의 꼭대기가 있을 수 있는 이유는 무엇일까? 조화 진동자가 그 예를 보여 준다. 소멸 연산자는 바닥 상태에 도달하면 0을 준다. 여기서는 올림 연산자가 꼭대기 상태에 도달하면 0을 주어야 한다.
\begin{equation}
L_+\lvert m_{\max}\rangle=0.
\end{equation}
\(L_+\)가 고윳값을 높인다는 말은 그 연산자가 만들어 내는 벡터가 0이 아닐 때만 성립한다. 최고 가중치 상태는 올림 과정이 끝나는 바로 그 예외적인 상태다.

회전 불변성은 아래쪽 끝점을 정해 준다. 각운동량이 \(+z\) 방향을 최대로 향할 수 있다면, 같은 물체를 회전시켜 \(-z\) 방향을 최대로 향하게 할 수도 있다. 따라서
\begin{equation}
m_{\min}=-m_{\max}.
\end{equation}
스펙트럼은 0을 중심으로 대칭이고 간격은 한 단위다. 가능한 계열은 두 가지다.

첫 번째 계열에는 0이 포함된다.\begin{equation}
\begin{aligned}
m={}&-m_{\max},-m_{\max}+1,\ldots,-1,0,\\
&1,\ldots,m_{\max}-1,m_{\max}.
\end{aligned}
\end{equation}그러면 \(m_{\max}\)는 정수이고 상태의 수는 다음과 같다.\begin{equation}
2m_{\max}+1.
\end{equation}최댓값은 입자의 특성을 나타내며 스핀이라 부른다. 통상적인 표기에서는\begin{equation}
s\equiv m_{\max},
\qquad
\lvert s,m\rangle\equiv\lvert m\rangle
\end{equation}고정된 스핀 다중항 안에서 이렇게 쓴다.\footnote{\raggedright 강의 영상 프레임 \nolinkurl{lecture_07_figure_03.png}, 01:01:08.\par}

\begin{figure}[htbp]
\centering
\includegraphics[width=0.70\textwidth]{lecture_07_figure_03.png}
\caption{최고 가중치 조건에서 올림 과정이 끝나며, 허용되는 \(m\) 값은 간격이 한 단위인 유한한 사다리를 이룬다.}
\lectureframe{01:01:08}
\end{figure}

스핀이 0인 입자에는 \(m=0\)만 있다. 스핀이 1인 입자는 다음과 같다.\[
m=1,0,-1,
\]스핀이 2인 입자는 성분 값이 다섯 개다. 힉스 보손은 스핀 0의 예다. 광자와 \(Z\) 보손은 스핀 1이다. 존재한다면 중력자는 스핀 2다. 복합 물체는 스핀 3이나 그보다 훨씬 높은 스핀도 가질 수 있다. 원자핵도 그럴 수 있고, 농구공이나 은하처럼 충분히 거시적인 회전 물체는 임의로 큰 각운동량을 가질 수 있다. 보통 기본 입자 목록에는 기본 스핀 3 입자가 없다.

대칭이고 간격이 한 단위인 다른 계열에는 0이 포함되지 않는다. 이 계열은 다음 값에서 시작한다.
\[
m=\frac12,-\frac12
\]
그리고 한 단위 간격으로 이어진다.
\begin{equation}
\begin{aligned}
m={}&-m_{\max},-m_{\max}+1,\ldots,-\frac32,-\frac12,\\
&\frac12,\frac32,\ldots,m_{\max}.
\end{aligned}
\end{equation}
여기서 \(m_{\max}\)는 반정수다. 가장 간단한 경우는 스핀 \(\tfrac12\)이고, 그다음은 스핀 \(\tfrac32\), 스핀 \(\tfrac52\) 등이다. 따라서 각운동량 다중항은 정수 스핀 계열과 반정수 스핀 계열로 나뉜다.

\begin{classroomqa}\lecturetimestamp{01:03:46}\audiencequestion{반정수 스핀 다중항에도 최댓값이 있나요?}\lecturerresponse{네. 이 사다리에도 최고값이 있다. 스핀이 \(\tfrac32\)인 입자의 최댓값은 \(m_{\max}=\tfrac32\)이고, 스핀이 \(\tfrac52\)인 입자의 최댓값은 \(m_{\max}=\tfrac52\)다. 차이는 최댓값과 다중항 안의 나머지 모든 값이 반정수라는 점뿐이다.}\end{classroomqa}

\begin{classroomqa}\lecturetimestamp{01:04:04}\audiencequestion{스핀 3 이상의 기본 입자는 없나요?}\lecturerresponse{여기서 다루는 상대론적 양자 틀에서 통상적인 목록에 포함되는 기본 입자의 스핀은 \(0,\tfrac12,1,\tfrac32,\)와 \(2\)다. 복합 물체는 이 범위에 제한되지 않으며 스핀 3이나 그보다 훨씬 큰 스핀을 가질 수 있다.}\end{classroomqa}

\begin{classroomqa}\lecturetimestamp{01:04:50}\audiencequestion{사다리에서 음수 값은 음의 스핀을 뜻하나요?}\lecturerresponse{아니다. 음의 \(m\)는 성분 값이다. 선택한 축을 기준으로 각운동량이 반대 방향을 향한다는 뜻이다. 입자의 스핀은 음이 아닌 최댓값 \(s=m_{\max}\)다.}\end{classroomqa}

\begin{classroomqa}\lecturetimestamp{01:05:10}\audiencequestion{성분 값이 불연속이라면 그 사이의 모든 방향은 어떻게 되나요?}\lecturerresponse{가능한 공간 방향이 둘이나 셋뿐인 것은 아니다. 중간 방향은 \(m\) 기저 상태들의 양자 중첩으로 나타낸다. 아래의 스핀 \(\tfrac12\) 구성에서 이를 명시적으로 볼 것이다.}\end{classroomqa}

\begin{classroomqa}\lecturetimestamp{01:05:42}\audiencequestion{코르크따개처럼 나아가는 운동이나 원편광은 어떤가요?}\lecturerresponse{원편광 전자기파는 원편광 광자들로 이루어지며, 이 광자들은 전파 축을 따라 각운동량을 지닌다. ``코르크따개''라는 말은 회전과 전진 운동을 함께 나타내므로, 정지 상태로 고려하는 물체를 설명하기에는 적절하지 않다. 움직이는 입자에서 관련된 물리량은 나선도다. 이는 운동 방향에 대한 스핀의 방향을 나타낸다.}\end{classroomqa}

\begin{classroomqa}\lecturetimestamp{01:07:10}\audiencequestion{스핀 \(\tfrac32\)인 기본 입자 가운데 알려진 것이 있나요?}\lecturerresponse{스핀 \(\tfrac32\)인 기본 입자는 직접 확인된 바 없다. 초대칭 이론은 중력자의 스핀 \(\tfrac32\) 짝입자인 중력자지노를 가정하지만, 이 입자도 직접 검출되지 않았다.}\end{classroomqa}

\begin{classroomqa}\lecturetimestamp{01:08:01}\audiencequestion{각운동량 연산자는 원래 파동함수나 켓 공간에 더 이상 작용하지 않나요?}\lecturerresponse{원래 스핀 공간은 아직 지정하지 않았다. 지금 대수를 사용해 그 공간을 정하고 있다. 스핀이 \(s\)로 고정되면 기저 벡터는 허용되는 값으로 표기한다.\[
\lvert s,m\rangle,\qquad m=-s,-s+1,\ldots,s.
\]이 고유 스핀 공간은 궤도 각운동량이 \(\mathbf r\times\mathbf p\)로 작용하는 보통의 공간 파동함수 공간과 구별되지만, 두 공간은 같은 각운동량 대수를 구현한다.}\end{classroomqa}

\section{각운동량의 크기}

\begin{classroomqa}\lecturetimestamp{01:09:32}\audiencequestion{각운동량 벡터의 크기는 관측 가능한가요?}\lecturerresponse{네. 그 제곱을 사용하는 것이 가장 편리하다.\[
L^2=L_x^2+L_y^2+L_z^2.
\]최고 가중치 상태에서 그 값은 단순히 \(s^2\hbar^2\)가 아니다. \(L_z\)가 정확히 정해져 최댓값을 갖더라도, 서로 교환하지 않는 가로 방향 성분 둘 다를 동시에 정확히 0으로 만들 수 없으므로 추가적인 양자 기여가 생긴다.}\end{classroomqa}

그 기여를 계산하려면 켓에 올림 연산자가 먼저 작용하는 순서로 사다리 연산자를 곱하자.\begin{align}
L_-L_+
&=(L_x-iL_y)(L_x+iL_y) \\
&=L_x^2+L_y^2+iL_xL_y-iL_yL_x \\
&=L_x^2+L_y^2+i[L_x,L_y] \\
&=L_x^2+L_y^2-\hbar L_z.
\end{align}마지막 단계에서는 \([L_x,L_y]=i\hbar L_z\)을 사용했다. 따라서\begin{equation}
L_x^2+L_y^2=L_-L_++\hbar L_z,
\end{equation}이고, 그러므로\begin{equation}
L^2=L_-L_++L_z^2+\hbar L_z.
\end{equation}

이 항등식을 최고 가중치 상태 \(\lvert s,s\rangle=\lvert m_{\max}\rangle\)에 적용하면 다음과 같다.
\begin{align}
L^2\lvert s,s\rangle
&=\left(L_-L_++L_z^2+\hbar L_z\right)\lvert s,s\rangle \\
&=L_-L_+\lvert s,s\rangle
+s^2\hbar^2\lvert s,s\rangle
+s\hbar^2\lvert s,s\rangle.
\end{align}
첫째 항은 \(L_+\lvert s,s\rangle=0\)이므로 사라진다. 따라서
\begin{equation}
L^2\lvert s,s\rangle
=s(s+1)\hbar^2\lvert s,s\rangle.
\end{equation}
추가 항 \(+s\)는 피할 수 없는 횡방향 요동에 따른 양자 보정이다. 고전적으로는 \(s^2\hbar^2\)만 예상했을 것이다.

처음 몇 가지 값을 보면 차이가 분명해진다.\begin{align}
s=0 &: \qquad L^2=0,\\
s=1 &: \qquad L^2=2\hbar^2,\\
s=\frac12 &: \qquad L^2=\frac34\hbar^2.
\end{align}\(s=1000\)일 때 무차원 인자는 \(1000\times1001\)이며, 큰 각운동량의 규모를 기준으로 보면 \(1000^2\)와 거의 다르지 않다. 일반적으로는\begin{equation}
s(s+1)=s^2\left(1+\frac1s\right),
\end{equation}이므로, \(s\)가 매우 클 때 상대적인 양자 보정은 무시할 수 있을 만큼 작아진다. 이것이 고전적인 \(s^2\) 결과가 나타나는 방식이다.

고정된 스핀 다중항의 모든 \(m\) 상태는 \(L^2\)의 고윳값이 같다. \(m\)를 바꾸면 선택한 축을 따른 성분만 바뀌고 전체 크기는 바뀌지 않는다. 이는 각운동량 벡터의 길이는 그대로 둔 채 방향을 회전하는 것에 해당하는 양자적 표현이다.

\begin{classroomqa}\lecturetimestamp{01:16:43}\audiencequestion{\(L^2\)와 \(z\) 성분을 동시에 알 수 있나요?}\lecturerresponse{네. 전체 제곱은 각 각운동량 성분과 모두 교환한다.\[
[L^2,L_i]=0,\qquad i=x,y,z.
\]따라서 서로 다른 두 성분은 동시에 지정할 수 없지만, \(L^2\)와 \(L_z\)처럼 선택한 한 성분은 함께 지정할 수 있다. 예를 들어 모든 전자는 방향과 관계없이 \(L^2=\tfrac34\hbar^2\)를 갖는다.}\end{classroomqa}

\begin{classroomqa}
\lecturetimestamp{01:17:41}
\audiencequestion{스핀 일이면 \(L^2\)은 0인가, 아니면 스핀 영일 때만 0인가?}
\lecturerresponse{\(L^2=0\)인 경우는 스핀 영뿐이다. 스핀 일에서는
\[
L^2=1(1+1)\hbar^2=2\hbar^2,
\]
스핀 삼에서는
\[
L^2=3(3+1)\hbar^2=12\hbar^2.
\]}
\end{classroomqa}

\begin{classroomqa}\lecturetimestamp{01:18:13}\audiencequestion{다중항에서 모든 \(m\) 값의 \(L^2\) 고윳값이 같은 이유는 무엇인가요?}\lecturerresponse{최고 가중치 상태에서 시작해 아래로 내려가 보자. \(L^2\)는 모든 성분과 교환하므로 \(L_-\)와도 교환한다.\[
[L^2,L_-]=0.
\]만약\[
L^2\lvert s,s\rangle=s(s+1)\hbar^2\lvert s,s\rangle,
\]라면,\begin{align}
L^2\bigl(L_-\lvert s,s\rangle\bigr)
&=L_-L^2\lvert s,s\rangle \\
&=s(s+1)\hbar^2\,L_-\lvert s,s\rangle.
\end{align}가 성립한다. 따라서 바로 아래 상태도 \(L^2\)의 고윳값이 같다. 같은 논리를 반복하면 그 고윳값이 사다리의 끝까지 유지된다.}\end{classroomqa}

따라서 불연속적인 \(m\) 값이 존재한다고 해서 공간 방향도 불연속적인 일부만 가능하다는 뜻은 아니다. 회전 대칭성이 있으므로 회전된 방향에 해당하는 상태들도 여전히 존재해야 한다. 스핀 \(\tfrac12\)에서는 그런 방향이 두 \(L_z\) 기저 상태의 중첩으로 나타난다.

\section{스핀\texorpdfstring{\(\tfrac12\)}{12} 상태와 방향}

회전 대칭성에 따라 불연속적인 \(L_z\) 고유상태만 있는 것으로는 충분하지 않다. 선택한 축에 대한 성분을 측정하면 허용된 불연속 값만 나오더라도 스핀은 어떤 축을 따라서도 향할 수 있다. \(z\)축 고유상태 사이의 방향은 중첩으로 표현된다.

스핀-\(\tfrac12\) 입자에서는 스핀 자유도만 살펴보자. 전자에는 궤도 상태도 있으므로, 전자의 전체 상태가 두 개뿐인 것은 결코 아니다. 스핀 공간만 보면 차원이 2다. \(z\)축을 기준으로 다음 기저를 택하자.
\begin{equation}
L_z\lvert+\rangle=\frac{\hbar}{2}\lvert+\rangle,
\qquad
L_z\lvert-\rangle=-\frac{\hbar}{2}\lvert-\rangle.
\end{equation}
일반적인 스핀 상태는 다음과 같다.
\begin{equation}
\lvert\psi\rangle
=
\alpha\lvert+\rangle+\beta\lvert-\rangle,
\end{equation}
여기서 \(\alpha\)와 \(\beta\)는 복소수다. \(L_z\)를 측정하면 확률 \(|\alpha|^2\)로 \(+\hbar/2\)을, 확률 \(|\beta|^2\)로 \(-\hbar/2\)을 얻는다. 따라서 정규화 조건은 다음과 같다.
\begin{equation}
|\alpha|^2+|\beta|^2=1.
\end{equation}

공통 위상은 물리적 효과가 없다.
\begin{equation}
(\alpha,\beta)\longrightarrow
e^{i\chi}(\alpha,\beta).
\end{equation}
두 복소 진폭은 처음에 네 개의 실수 매개변수로 표현된다. 정규화 조건으로 하나가 정해지고, 물리적으로 무의미한 전체 위상으로 또 하나가 정해지므로 물리적 실수 매개변수는 둘 남는다. 이는 방향을 극각과 방위각으로 나타내는 데 필요한 수와 정확히 같다. \(\alpha\)와 \(\beta\)를 바꾸면 스핀 방향이 회전한다.

측정된 성분과 그 기댓값은 서로 다른 것이다. 기댓값은 앙상블 평균이다. 같은 상태로 준비한 전자를 많이 두고 각 전자에 대해 \(L_z\)를 한 번씩 측정한 뒤, 결과를 모두 더해 실험 횟수로 나누자. 각 측정 결과는 여전히 \(+\hbar/2\) 또는 \(-\hbar/2\)뿐이지만, 평균은 그 두 값 사이의 임의의 값을 취할 수 있다.
\begin{equation}
\langle L_z\rangle
=
\frac{\hbar}{2}\bigl(|\alpha|^2-|\beta|^2\bigr).
\end{equation}

스핀 업과 스핀 다운의 진폭이 같으면 스핀은 \(x\)축의 두 방향 가운데 한쪽을 향한다. 정규화된 상태의 표현은\footnote{계수의 비가 방향을 결정한다. \(1/\sqrt2\) 인자는 \( |\alpha|^2+|\beta|^2=1\)이 요구하는 정규화를 제공한다. 상대 위상 \(+i\)과 \(-i\)은 서로 반대인 \(y\)축 방향을 나타낸다. 어느 부호를 양의 \(y\) 방향으로 정할지는 축의 규약에 달려 있다.}\begin{equation}
\lvert x;+\rangle
=
\frac{\lvert+\rangle+\lvert-\rangle}{\sqrt2},
\qquad
\lvert x;-\rangle
=
\frac{\lvert+\rangle-\lvert-\rangle}{\sqrt2}.
\end{equation}두 상태 중 어느 상태에서든 \(z\)축 방향으로 측정하면 가능한 두 결과가 같은 확률로 나오므로\begin{equation}
\langle L_z\rangle
=
\frac12\left(\frac{\hbar}{2}\right)
+
\frac12\left(-\frac{\hbar}{2}\right)
=0.
\end{equation}\(x\)축 방향에서는 상대 부호가 플러스인 상태와 마이너스인 상태가 서로 반대 방향을 나타낸다. 마찬가지로 \begin{equation}
\frac{\lvert+\rangle+i\lvert-\rangle}{\sqrt2},
\qquad
\frac{\lvert+\rangle-i\lvert-\rangle}{\sqrt2}
\end{equation}은 서로 반대인 \(y\)축 방향을 나타낸다. \(\alpha\)과 \(\beta\)의 일반적인 정규화 값은 임의의 축 방향을 나타낸다. 스핀이 특정 축을 향한다고 말하는 것은 그 축 방향 성분을 측정했을 때 확정된 결과가 나온다는 뜻이다. 다른 축을 기준으로 보면 같은 상태는 일반적으로 중첩 상태이다.

개별 측정 결과는 여전히 양자화되어 있다. 연속적으로 변하며 고전적 벡터의 익숙한 기하학적 성질을 재현하는 것은 앙상블 평균이다.

\section{스핀, 통계, 복합 입자}

질량, 전하, 스핀은 입자의 특성을 나타내는 주요 물리량이다. 스핀은 통계와도 연관된다. 스핀-통계 관계는 상대론적 양자역학의 정리지만, 여기서는 증명하지 않겠다.
\begin{equation}
\begin{aligned}
\text{half-integer spin} &\quad\Longleftrightarrow\quad \text{fermion},\\
\text{integer spin} &\quad\Longleftrightarrow\quad \text{boson}.
\end{aligned}
\end{equation}
동일한 페르미온 두 개는 완전하게 같은 양자 상태를 차지할 수 없다. 동일한 보손들은 같은 상태를 차지할 수 있으며, 많은 보손이 한 상태에 있으면 집단적으로 고전파처럼 거동할 수 있다. 따라서 스핀 \(\tfrac12,\tfrac32,\tfrac52,\ldots\)인 입자는 페르미온이고, 스핀 \(0,1,2,\ldots\)인 입자는 보손이다.

이 규칙은 복합 입자에도 바로 적용된다. 반정수 각운동량 두 개를 더하면 정수 각운동량이 된다. 더 일반적으로 페르미온 구성 요소가 짝수 개이면 복합 입자는 보손이 되고, 홀수 개이면 페르미온이 된다. 보손 구성 요소를 더해도 이 홀짝 규칙은 바뀌지 않는다.

보통의 수소 원자에는 스핀-\(\tfrac12\) 양성자와 스핀-\(\tfrac12\) 전자가 하나씩 있다. 따라서 페르미온이 짝수 개이며, 전체적으로 보손이다. 중수소핵에는 양성자와 중성자가 하나씩 들어 있으며, 둘 다 페르미온이므로 이 역시 보손이다. 중수소 원자에는 중수소핵에 전자가 하나 더해진다. 여기서 사용한 설명 수준에서는 페르미온 구성 요소가 세 개이므로 중수소 원자는 페르미온이다.

입자와 반입자는 질량과 스핀이 같고 전하가 서로 반대다. 포지트로늄은 전자와 양전자가 서로 결합한 상태다. 페르미온 두 개로 이루어져 있으므로 보손이다.

\begin{classroomqa}
\lecturetimestamp{01:33:08}
\audiencequestion{보손은 물질을 통과해 지나가며, 초전도는 짝을 이룬 전자들이 보손처럼 거동하는 것과 관련이 있는가?}
\lecturerresponse{관련된 성질은 물질을 통과하는 능력 자체가 아니라 보손이 같은 상태를 차지할 수 있다는 점이다. 더 간단한 예는 초유체 헬륨이다. 보통의 헬륨-\(4\) 원자는 양성자 두 개, 중성자 두 개, 전자 두 개를 포함한다. 페르미온 구성 요소가 짝수 개이므로 이 원자는 보손이다. 많은 헬륨 원자가 한 양자 상태를 차지하고 집단적으로 움직일 수 있다. 그러면 초유체를 많은 원자를 같은 상태에 쌓아 형성한 고전파로 근사적으로 그려 볼 수 있다. 이 그림은 정확하지는 않지만, 여기서 필요한 요점을 잘 보여 준다. 금속 안의 전자는 페르미온이므로 다른 전자와 완전하게 같은 상태를 공유할 수 없다. 그러나 전자쌍은 보손인 복합 입자이며, 이것이 초전도와의 관련성이다. 여기서 역설이 생긴다. 구성 전자들이 같은 상태를 차지할 수 없다면, 복합 전자쌍 여러 개가 어떻게 한 상태를 차지할 수 있을까? 이를 해명하려면 복합 상태를 더 자세히 다뤄야 하므로 여기서는 미룬다.}
\end{classroomqa}

\begin{classroomqa}
\lecturetimestamp{01:35:46}
\audiencequestion{양성자와 중성자는 복합 입자인데 왜 반정수 스핀을 갖는가?}
\lecturerresponse{각 입자는 스핀-\(\tfrac12\) 쿼크가 홀수 개 모여 만들어진다. 페르미온 구성 요소 세 개는 페르미온인 복합 입자를 이루므로, 양성자와 중성자는 반정수 스핀을 가질 수 있다.}
\end{classroomqa}

\section{파울리 배타 원리와 원자 오비탈}

주기율표를 보면 파울리 배타 원리가 왜 필요한지 바로 알 수 있다. 이 원리는 전자 두 개가 같은 궤도 운동을 할 수 없다는 뜻이 아니다. 두 전자가 완전하게 같은 양자 상태를 차지할 수 없다는 뜻이다.

이 구분은 가장 단순한 원자에서도 중요하다. 수소 원자에서 시작해 핵전하를 늘리고 전자를 하나 더해 보자. 두 전자는 같은 공간 오비탈을 차지할 수 있지만, 스핀 상태는 달라야 한다. 선택한 축을 기준으로 한 전자는 스핀 업, 다른 전자는 스핀 다운일 수 있다. 오비탈 하나에 전자 두 개가 있으면 두 스핀은 스핀 단일항 상태를 이루므로 전체 스핀은 상쇄된다. 이용할 수 있는 독립적인 스핀 상태가 둘뿐이므로 공간 오비탈 하나에는 전자를 두 개까지만 넣을 수 있다.

\begin{classroomqa}
\lecturetimestamp{01:37:53}
\audiencequestion{오비탈 하나에 있는 두 전자의 반대 스핀을 맞물린 톱니바퀴로 그려 볼 수 있는가?}
\lecturerresponse{대략적인 그림으로는 그렇다. 맞물린 톱니바퀴는 서로 반대 방향으로 돈다. 다만 이것은 스핀 방향이 서로 반대임을 비유한 것일 뿐, 전자를 실제 기계 장치로 나타낸 모형은 아니다.}
\end{classroomqa}

``오비탈''이라는 말은 서로 다른 두 수준에서 쓰일 수 있다. 넓은 의미에서 궤도 운동 상태란 스핀을 무시하고 슈뢰딩거 방정식을 풀어 얻는 상태다. 이는 전자가 어느 껍질에 있는지만 나타내는 것보다 많은 정보를 담고 있다.

\begin{classroomqa}
\lecturetimestamp{01:38:09}
\audiencequestion{오비탈은 전자가 어느 껍질에 있는지만 나타내는 말보다 더 넓은 의미인가?}
\lecturerresponse{그렇다. 여기서 궤도 운동은 전자의 스핀을 무시하고 얻는 슈뢰딩거 상태를 뜻한다. 여기에는 전자의 반지름 방향 정보와 핵으로부터의 거리를 대략 정하는 주양자수 \(n\)가 들어가며, 궤도 각운동량 양자수도 포함된다. 화학자들이 사용하는 분광학 표기에서 \(S\), \(P\), \(D\), \(F\)는 각각 궤도 각운동량 \(l=0,1,2,3\)에 대응한다.}
\end{classroomqa}

따라서 주양자수와 궤도 각운동량은 운동의 서로 다른 측면을 나타낸다. \(S,P,D,F\) 표기는 \(l\)의 값에 해당하고, \(n\)는 별도의 반지름 방향 정보를 나타낸다. 파울리 원리에 따르면 스핀 양자수를 포함한 모든 양자수가 같은 전자는 둘 이상 있을 수 없다.

\begin{classroomqa}
\lecturetimestamp{01:40:42}
\audiencequestion{양자수로 원자의 각 껍질에 들어갈 수 있는 전자 수를 정할 수 있는가?}
\lecturerresponse{그렇다. 양자수로 가능한 최대 전자 수를 정할 수 있지만, 특정 원자의 해당 껍질에 실제로 전자가 몇 개 있는지까지 정해 주지는 않는다. 일반적인 껍질별 전자 수 계산은 여기서 다루지 않겠다.}
\end{classroomqa}

궤도 각운동량 \(l\)이 정해지면 성분 양자수는 다음과 같이
\begin{equation}
2l+1
\end{equation}
개의 가능한 값을 갖는다. 이는 앞서 스핀에서 다룬 각운동량의 상태 수와 같다. \(S\) 상태는 \(l=0\)이므로 궤도 상태가 하나다. \(P\) 상태는 \(l=1\)이므로 궤도 상태가 세 개다. \(D\) 상태는 \(l=2\)이므로 다섯 개다.

각 공간 오비탈에 반대 스핀 상태의 전자를 하나씩 넣을 수 있으므로, 스핀을 고려하면 각 궤도 상태의 수가 두 배가 된다.

\begin{classroomqa}
\lecturetimestamp{01:41:54}
\audiencequestion{\(S\) 상태에는 전자 두 개가 들어가고, \(P\) 상태 세 개에는 전자가 여섯 개 더 들어가는가?}
\lecturerresponse{그렇다. \(S\) 오비탈 상태 하나에는 전자가 최대 두 개 들어간다. \(P\) 오비탈 상태 세 개에는 각각 두 개씩 들어가므로 여섯 개가 더 들어간다. 따라서 여기서 다루는 \(S\) 상태와 \(P\) 상태에는 모두 \(2+6=8\)개의 전자가 들어간다.}
\end{classroomqa}

원자 오비탈을 구성하는 각운동량 구조는 약간의 차이는 있지만, 쿼크가 결합계 안에서 움직이며 서로를 공전할 때도 다시 나타난다.
